\clearpage
\fancyhead[L]{\small GAUSS AN BODE.}
\hypertarget{native-438-part-lower}{}
\begin{center}
{\large GAUSS AN BODE.}

Astronomisches Jahrbuch für 1823 Seite 228..230. Berlin 1820.
\end{center}

\hfill Göttingen, den 3. Sept. 1820.

\par\indent Ihrem Wunsche zufolge übersende ich Ihnen hier die Ephemeride für den Lauf der
\emph{Pallas} während ihrer Sichtbarkeit im Jahre 1821. Sie ist diesmal von Hrn. \textsc{von Staudt} berechnet, einem jungen Manne von ausgezeichneten Talenten, welcher sich
hier dem Studium der Mathematik und Astronomie widmet. Da die Berechnung derselben erst
heute fertig geworden ist, und Sie dieselbe sobald als möglich zu erhalten wünschten, so
kann ich die noch nicht vollendete genaue Vorausberechnung der am 19. Mai 1821
einfallenden Opposition nicht beifügen, weil sonst die Absendung erst einen Posttag
später geschehen könnte.

Von dem \textsc{Reichenbach}schen Meridiankreise habe ich in unsern gel. Anz. eine
ausführliche Beschreibung gegeben, und es wird daher unnöthig sein, hier etwas davon zu
wiederholen. Die Anzahl meiner vom 21. Febr. bis jetzt damit gemachten Beobachtungen mag
etwa 1200 betragen, von denen aber bis jetzt erst ein Theil ganz vollständig reducirt
ist. Ich schreibe Ihnen diesmal meine bis zum 15. Juli beobachteten
Sonnendeclinationen her (von den gleichzeitig am Mittagsfernrohr beobachteten
Rectascensionen ist erst ein Theil reducirt). Der Ort des Pols auf dem Instrument ist
dabei bereits nach der Gesammtheit der dazu dienlichen Beobachtungen auf das schärfste
bestimmt, und die Beobachtungen sind von der Refraction nach \textsc{Bessels} Tafel und
von der Parallaxe, die mittlere Horizontalparallaxe $=8''60$ angenommen, befreit. Auf
Flexion ist aber keine Rücksicht genommen. Nach meinen bisherigen Versuchen scheint
dieselbe fast ganz unmerklich zu sein, und wenn sie doch vielleicht \emph{einige
Zehntheile der Secunde} betragen sollte, so kann diese Grösse erst durch eine lange
Reihe von Beobachtungen ausgemittelt werden.


\Needspace{17\baselineskip}
\SourceLeaf{439}
\begin{center}
\emph{Beobachtete Declinationen der Sonne} 1820.

O obern \quad U untern Rand
\end{center}

\begin{center}
\begin{tabular}{ll|r r r l}
\multicolumn{6}{c}{Südlich} \\
Febr. & 22 & $10^\circ$ & $12'$ & $36''5$ & O \\
 & 27 & $8$ & $54$ & $23.0$ & U \\
 & 28 & $7$ & $59$ & $29.9$ & O \\
 & 29 & $8$ & $9$ & $11.3$ & U \\
März & 5 & $6$ & $14$ & $26.5$ & U \\
 & 11 & $3$ & $53$ & $5.3$ & U \\
 & 14 & $2$ & $11$ & $0.0$ & O \\
 & 19 & $0$ & $44$ & $37.2$ & U \\
\multicolumn{6}{c}{Nördlich} \\
April & 3 & $5$ & $38$ & $52.1$ & O \\
 & 4 & $6$ & $1$ & $48.2$ & O \\
\end{tabular}
\end{center}

\Needspace{15\baselineskip}
\begin{center}
\begin{tabular}{ll|r r r l}
\multicolumn{6}{c}{Nördlich} \\
April & 12 & $9^\circ$ & $0'$ & $52''2$ & O \\
 & 13 & $8$ & $50$ & $43.1$ & U \\
 & 16 & $9$ & $55$ & $14.7$ & U \\
 & 18 & $11$ & $9$ & $15.3$ & O \\
 & 21 & $12$ & $11$ & $13.2$ & O \\
 & 22 & $11$ & $59$ & $26.2$ & U \\
 & 23 & $12$ & $19$ & $29.7$ & O \\
 & 24 & $13$ & $11$ & $9.9$ & O \\
 & 26 & $13$ & $50$ & $10.3$ & O \\
 & 27 & $13$ & $37$ & $30.7$ & U \\
Mai & 6 & $16$ & $19$ & $13.5$ & U \\
\end{tabular}
\end{center}

\Needspace{15\baselineskip}
\begin{center}
\begin{tabular}{ll|r r r l}
\multicolumn{6}{c}{Nördlich} \\
Mai & 11 & $18^\circ$ & $11'$ & $25''9$ & O \\
 & 19 & $20$ & $4$ & $23.5$ & O \\
 & 27 & $21$ & $4$ & $11.3$ & U \\
 & 31 & $21$ & $41$ & $11.8$ & U \\
Juni & 6 & $22$ & $25$ & $19.1$ & U \\
 & 20 & $23$ & $43$ & $27.5$ & O \\
 & 25 & $23$ & $50$ & $30.6$ & O \\
 & 27 & $23$ & $4$ & $58.1$ & U \\
Juli & 11 & $22$ & $23$ & $0.0$ & O \\
 & 12 & $22$ & $14$ & $51.1$ & O \\
 & 15 & $21$ & $16$ & $29.3$ & U \\
\end{tabular}
\end{center}

Noch theile ich Ihnen die Resultate meiner Beobachtungen von 20 Zenithalsternen mit,
d. i. die auf den Anfang von 1820 reducirten mittlern Zenithdistanzen. Zur Reduction
sind dieselben Elemente gebraucht, die \textsc{Bessels} Tafeln für den Nordstern zum
Grunde liegen. Es sind dies dieselben Sterne, welche Herr Prof. \textsc{Schumacher} mit
dem \textsc{Ramsden}schen Zenithsector 1819 in Lauenburg und in d. J. in Skagen
beobachtet hat, und woraus die Krümmung der verschiedenen Stücke des Meridianbogens von
Skagen bis Göttingen werden abgeleitet werden. Ich kann diese aber noch nicht angeben,
da mir die \emph{Endresultate} der \textsc{Schumacher}schen Beobachtungen noch nicht
bekannt sind. Ich bemerke nur noch, dass ich jeden Stern \emph{wenigstens} dreimal in
der einen und dreimal in der andern Lage des Kreises beobachtet habe. Bei der Reduction
auf den Anfang von 1820 ist auf die eigne Bewegung der Sterne keine Rücksicht genommen,
diese kann aber mit grösster Genauigkeit aus der Vergleichung meiner Bestimmungen mit
den \textsc{Brandley}schen
abgeleitet werden, wozu die Polhöhe der hiesigen Sternwarte $=51^\circ\,31'\,48''7$
anzunehmen ist. Das Mittel meiner Beobachtungen fällt in die ersten Tage des August.

\Needspace{27\baselineskip}
\begin{center}
\begin{tabular}{l|l|l}
Sterne. & M. Z. D. 1820. & Coll. Fehler. \\
\hline
$\eta$ Cephei & $9^\circ\,36'\,44''28$ N. & $2'\,15''50$ \\
\emph{Piazzi XX.} 222. & $8\;\;17\;\;7.34$ & $16.65$ \\
$\iota$ Draconis & $8\;\;4\;\;11.09$ & $16.78$ \\
47 Draconis & $7\;\;38\;\;26.35$ & $17.19$ \\
Cephei 2 Hev. & $6\;\;50\;\;16.06$ & $16.96$ \\
48 Draconis & $6\;\;2\;\;55.95$ & $17.36$ \\
53 Draconis & $5\;\;1\;\;32.34$ & $17.06$ \\
33 Cygni & $4\;\;29\;\;25.71$ & $18.35$ \\
49 Draconis & $3\;\;52\;\;25.79$ & $17.18$ \\
46 Draconis & $3\;\;49\;\;49.51$ & $16.62$ \\
\emph{Piazzi XX.} 391. & $2\;\;17\;\;59.29$ & $17.64$ \\
51 Draconis & $1\;\;35\;\;38.56$ & $16.96$ \\
$\kappa$ Cygni & $1\;\;30\;\;37.18$ & $17.20$ \\
\emph{Piazzi XXI.} 32. & $1\;\;18\;\;6.93$ & $16.27$ \\
20 Cygni & $1\;\;0\;\;18.21$ & $17.54$ \\
$\iota$ Cygni & $0\;\;10\;\;46.28$ S. & $17.06$ \\
$c'$ Cygni Praec. & $1\;\;25\;\;6.27$ & $18.03$ \\
$c'$ Cygni Sequ. & $1\;\;25\;\;33.33$ & $16.37$ \\
$\theta$ Cygni & $1\;\;43\;\;15.97$ & $17.79$ \\
$\omega'$ Cygni & $2\;\;44\;\;22.59$ & $17.71$ \\
\multicolumn{2}{r|}{Mittlerer Werth des C. F. . .} & $2'\,17''22$ \\
\end{tabular}
\end{center}

Sie sehen, dass nur bei zwei Sternen die Bestimmung des Collimationsfehlers über $1''$
vom mittlern Werthe abweicht.
